\chapter{Traitement numérique des mesures expérimentales}

Ce chapitre est utile lorsque nous avons obtenu des mesures expérimentales et
voulons les traiter. Les techniques numériques utilisées pour le traitement de
ces données portent essentiellement sur l'interpolation et l'approximation des
fonctions.

\section{Généralités}

\subsection{Position du problème}

L'interpolation et l'approximation des fonctions sont des techniques
numériques très utilisées en sciences expérimentales. Le problème est le
suivant :

Nous avons une grandeur $y$ qui dépend d'une variable $x$. Nous connaissons
les valeurs $y_i$ de cette grandeur en un certain nombre de valeurs discrètes
ou points discrets $x_1, x_2, \dots, x_n$. Nous voulons approximer $y$ par une
fonction analytique $y = f(x)$ telle que
\[
  y_1 = f(x_1),\quad y_2 = f(x_2),\quad \dots,\quad y_n = f(x_n).
\]
L'expression de $f(x)$ pourra ainsi permettre de déterminer les valeurs de la
fonction en des points autres que ceux de la suite discrète. On parle alors de
l'art de lire entre les lignes d'une table ou d'\emph{interpolation}. Ces
points où on cherche les nouvelles valeurs de $y$ peuvent être inaccessibles
expérimentalement (soit parce qu'ils sont éloignés de la plage de mesure de
l'appareil utilisé, soit parce que la précision de l'appareil de mesure ne
permet pas d'atteindre ces points).

Un tel problème n'admet pas de solution univoquement déterminée. Notons aussi
que l'expression analytique écrite de $f(x)$ ne nous intéresse pas
particulièrement puisqu'elle peut être gérée par l'ordinateur pour lire entre
les lignes d'une table. Mais on peut avoir besoin d'une telle expression pour
les besoins d'analyse des courbes.

Le problème peut se poser dans des cas où la grandeur $y$ dépend d'une seule
variable $x$, de deux ou de trois variables.

\subsection{Type d'interpolation}

La fonction $f(x)$ peut avoir plusieurs formes. Lorsqu'elle est un polynôme,
on parle d'interpolation \textbf{polynomiale}. Lorsque $f(x)$ est une fonction
trigonométrique, on parle d'interpolation \textbf{trigonométrique}. On peut
aussi avoir une interpolation par les fonctions exponentielles, des polynômes
de Legendre, les fonctions de Bessel, etc. Chacune de ces formes est choisie
en fonction des prédictions dictées par la nature du problème et ce que l'on
attend.

Cependant en pratique, l'on choisit souvent la forme la plus simple pour
$f(x)$, notamment la forme polynomiale. Dans le cas où l'on sait que la
fonction est périodique, il est judicieux d'utiliser l'interpolation
trigonométrique. La justification de l'un ou de l'autre type d'interpolation
repose sur 2 théorèmes présentés par \textbf{Weierstrass} (en 1885).

\paragraph{\underline{Théorème 1}} Toute fonction continue dans un intervalle
$[a,b]$ peut être représentée dans cet intervalle et pour chaque degré de
précision par un polynôme
\[
  f(x) = \sum_{n=0}^{N} a_n x^n .
\]

\paragraph{\underline{Théorème 2}} Toute fonction continue de période $2\pi$
peut être représentée par une série trigonométrique limitée de la forme
\[
  f(x) = \sum_{n=0}^{N} \bigl(a_n\cos nx + b_n\sin nx\bigr).
\]
Les coefficients $a_n$ et $b_n$ sont déterminés par la connaissance des
valeurs discrètes de la fonction sur l'intervalle $[a,b]$.

\section{Interpolation polynomiale}

Un polynôme de degré $n$ est défini par ses $(n+1)$ coefficients. Ainsi un
polynôme d'interpolation de degré $n$ sera complètement déterminé par la
connaissance de $(n+1)$ valeurs discrètes $y_i$. Il existe plusieurs formules
pour déterminer l'interpolation polynomiale.

\subsection{Formule d'interpolation de Lagrange}

L'idée de base de la formule de Lagrange est de déterminer des polynômes
$f_k(x)$ de degré $n$ qui prennent les valeurs 1 au point $x_k$ et 0 en tout
point différent de $x_k$
\[
  \Longrightarrow\quad
  f_k(x_{k'}) =
  \begin{cases}
    1 & \text{si } k' = k\\
    0 & \text{si } k' \neq k
  \end{cases}
\]
Le polynôme d'interpolation de Lagrange aura alors la forme
\[
  f(x) = y_1f_1(x) + y_2f_2(x) + \dots + y_{n+1}f_{n+1}(x)
  \quad\Longrightarrow\quad
  f(x) = \sum_{k=1}^{n+1} y_kf_k(x).
\]

\paragraph{Détermination de $f_k(x)$}

Soit par exemple un ensemble de deux points $(x_1, y_1)$ et $(x_2, y_2)$. Les
polynômes $f_k$ sont de degré 1, c'est-à-dire
\[
  f_1(x) = a_1 + b_1x\,; \qquad f_2(x) = a_2 + b_2x
\]
\[
  \left\{
  \begin{aligned}
    f_1(x_1) &= 1\\
    f_1(x_2) &= 0
  \end{aligned}
  \right.
  \quad\text{et}\quad
  \left\{
  \begin{aligned}
    f_2(x_1) &= 0\\
    f_2(x_2) &= 1
  \end{aligned}
  \right.
  \quad\Longrightarrow\quad
  f_1(x) = \frac{x - x_2}{x_1 - x_2}, \qquad
  f_2(x) = \frac{x - x_1}{x_2 - x_1}.
\]
Donc $y = y_1f_1(x) + y_2f_2(x)$.

Si on avait un ensemble de trois points $(x_1, y_1)$, $(x_2, y_2)$,
$(x_3, y_3)$, alors les polynômes seront de degré 2 et définis par :
\[
  f_1(x) = \frac{(x - x_2)(x - x_3)}{(x_1 - x_2)(x_1 - x_3)}, \quad
  f_2(x) = \frac{(x - x_1)(x - x_3)}{(x_2 - x_1)(x_2 - x_3)}, \quad
  f_3(x) = \frac{(x - x_1)(x - x_2)}{(x_3 - x_1)(x_3 - x_2)} .
\]
Dans ce cas, on a $y = y_1f_1(x) + y_2f_2(x) + y_3f_3(x)$. On parle alors
d'\emph{interpolation quadratique}.

En généralisant pour un ensemble de $n+1$ points, les polynômes de degré $n$
tels que
\[
  f_k(x_{k'}) =
  \begin{cases}
    1 & \text{si } k' = k\\
    0 & \text{si } k' \neq k
  \end{cases}
\]
sont définis par
\[
  f_k(x) = \frac{\displaystyle\prod_{\substack{i=1\\ i\neq k}}^{n+1}(x - x_i)}
                {\displaystyle\prod_{\substack{i=1\\ i\neq k}}^{n+1}(x_k - x_i)}
  \qquad\text{et}\qquad
  f(x) = \sum_{k=1}^{n+1} f_k(x)\,y_k .
\]
Le polynôme $f(x)$ est unique pour un ensemble de points $x_i$.

\paragraph{Exemple} On donne le tableau suivant :
\begin{center}
  \begin{tabular}{|c|c|c|c|}
    \hline
    $x$ & $-1$ & $0$ & $1$\\ \hline
    $y$ & $4$ & $1$ & $2$\\ \hline
  \end{tabular}
\end{center}
$(x_1 = -1,\ x_2 = 0,\ x_3 = 1)$. D'après la formule de Lagrange, on a :
\[
  \begin{aligned}
    f_1(x) &= \frac{(x - x_2)(x - x_3)}{(x_1 - x_2)(x_1 - x_3)} = \frac{x(x-1)}{2}, &
    f_2(x) &= \frac{(x - x_1)(x - x_3)}{(x_2 - x_1)(x_2 - x_3)} = -(x^2 - 1),\\
    f_3(x) &= \frac{(x - x_1)(x - x_2)}{(x_3 - x_1)(x_3 - x_2)} = \frac{(x+1)x}{2}
  \end{aligned}
\]
d'où
\[
  \begin{aligned}
    f(x) &= y_1f_1(x) + y_2f_2(x) + y_3f_3(x)\\
    f(x) &= 2x(x-1) - (x^2 - 1) + x(x+1).
  \end{aligned}
\]

À partir de l'expression de la fonction de Lagrange, on peut évaluer les
valeurs $y_p$ pour toute valeur $x_p$ comprise entre deux points $x_k$
(interpolation) ou en dehors des points $x_k$ (extrapolation). Pour ce faire,
on utilise l'algorithme suivant :

\begin{algo}
  \State 1- Entrer les données $(x_k, y_k, x_p)$ avec $k = 1, 2, \dots, n+1$
  \State 2- $y_p \leftarrow 0$
  \State 3- \textbf{Pour} $k$ allant de $1$ jusqu'à $n+1$, \textbf{faire}
  \State \hspace{2.2em} $y_s \leftarrow 1$
  \State \hspace{2.2em} \textbf{Pour} $i$ allant de $1$ jusqu'à $n+1$, \textbf{faire}
  \State \hspace{4.4em} \textbf{si} $i \neq k$
  \State \hspace{5.5em} $y_s \leftarrow y_s\,\dfrac{(x_p - x_i)}{(x_k - x_i)}$
  \State \hspace{2.2em} \textbf{Fin pour}
  \State \hspace{2.2em} $y_p \leftarrow y_p + y_k\,y_s$
  \State \hspace{1em} \textbf{Fin pour}
  \State 4- \textbf{Écrire} $x_p$ et $y_p$
  \State \textbf{Fin}
\end{algo}

\begin{exercice}
Traduire cet algorithme dans un des langages de programmation de votre choix
(Pascal, Fortran, C).
\end{exercice}

\subsection{Formule d'interpolation de Newton}

\subsubsection{Par les différences divisées}

On appelle différence divisée première des variables $x_i$ et $x_j$
correspondant aux valeurs $y_i$ et $y_j$ la quantité $\delta(x_i, x_j)$
\[
  \delta(x_i, x_j) = \frac{y_i - y_j}{x_i - x_j} .
\]
La différence divisée seconde est
\[
  \delta(x_i, x_j, x_k) = \frac{\delta(x_i, x_j) - \delta(x_i, x_k)}{x_j - x_k} .
\]
La différence divisée d'ordre $m$ est
\[
  \delta(x_1, \dots, x_m) =
  \frac{\delta(x_1, \dots, x_{m-1}) - \delta(x_1, \dots, x_{m-2}, x_m)}{x_{m-1} - x_m} .
\]
Le polynôme d'interpolation de Newton se met sous la forme
\[
  \begin{aligned}
    y = y_1 &+ \delta(x_1, x_2)(x - x_1) + \delta(x_1, x_2, x_3)(x - x_1)(x - x_2) + \dots\\
            &\dots + \delta(x_1, x_2, \dots, x_{k+1})(x - x_1)(x - x_2)\cdots(x - x_k) .
  \end{aligned}
\]

\subsubsection{Par les différences progressives et régressives}

La quantité $\Delta y_k = y_{k+1} - y_k$ est appelée différence première
progressive. La différence seconde progressive est
$\Delta^2 y_k = \Delta y_{k+1} - \Delta y_k$ et la différence progressive
d'ordre $\alpha$ est définie par
\[
  \Delta^{\alpha} y_k = \Delta^{\alpha-1} y_{k+1} - \Delta^{\alpha-1} y_k .
\]
On définit également la différence régressive d'ordre $\alpha$ de la manière
suivante
\[
  \nabla^{\alpha} y_k = \nabla^{\alpha-1} y_k - \nabla^{\alpha-1} y_{k-1}
  \qquad (\nabla y_k = y_k - y_{k-1}).
\]
Lorsque l'on a $x_{k+1} = x_k + h = x_1 + kh$, la formule d'interpolation par
les différences divisées se réduit à :
\[
  f(x) = y_1 + \frac{\Delta y_1}{h}(x - x_1)
             + \frac{\Delta^2 y_1}{h^2\,2!}(x - x_1)(x - x_2) + \dots
             + \frac{\Delta^{k} y_1}{h^{k}\,k!}(x - x_1)(x - x_2)\cdots(x - x_k) .
\]
Cette dernière formule est souvent appelée \textbf{formule de
Newton-Gregory}.

Si par contre on a $x_{k+1} = x_k - h = x_1 - kh$, alors on peut utiliser la
formule de Newton-Gregory régressive en remplaçant dans la formule précédente
$\Delta$ par $\nabla$.

\subsection{Erreurs d'interpolation polynomiale}

Dans la méthode de Lagrange, en remplaçant la valeur exacte $y_{\exp}$ (que
l'on aurait obtenue expérimentalement) par $y_p$, on commet une erreur de
troncature qui peut être déterminée par le théorème suivant :

\paragraph{\underline{Théorème}} Soit $f(x)$ le polynôme de degré $n$ qui
satisfait la condition $f(x_k) = y_k$ $(k = 1, \dots, n+1)$ et définie dans
l'intervalle $[a,b]$ qui contient les points $x_k$. Si $f$ est de classe
$C^{n+1}$ sur $[a,b]$ ($n+1$ fois dérivable), alors pour toute valeur
$x_p \in [a,b]$, il existe une valeur $\xi \in [a,b]$ telle que :
\[
  y_{\exp} - y_p = \Bigl[f^{(n+1)}(\xi)\big/(n+1)!\Bigr]
                   (x_p - x_1)(x_p - x_2)\cdots(x_p - x_{n+1}) .
\]
Cette différence représente l'\textbf{erreur de troncature}.

En supposant que les dérivées de la fonction $f$ sont uniformément bornées sur
$[a,b]$, l'erreur de troncature \textbf{décroît lorsque $n$ croît}.

\subsection{Interpolation Spline}

\subsubsection{Problème et propriétés}

L'interpolation Spline a originellement une formulation en mécanique. Elle est
utilisée (en méthodes des éléments finis) pour résoudre numériquement les
équations aux dérivées partielles. Son importance relève des inconvénients de
l'interpolation de Lagrange.

Le problème associé à l'interpolation Spline est le suivant :

Soit $x_1, x_2, \dots, x_n$ points distincts compris dans un intervalle
$[a,b]$ et soit $f_1, f_2, \dots, f_n$ $(y_i = f_i)$ les valeurs de la fonction
en ces points, l'interpolation Spline est une interpolation cubique qui passe
par les points $(x_i, f_i)$. C'est une fonction $S$ qui admet les propriétés
suivantes :
\begin{enumerate}[label=(\alph*)]
  \item $S(x_i) = f_i$,
  \item $S(x)$, $S'(x)$, $S''(x)$ sont continues sur $[a,b]$,
  \item dans chaque intervalle $[x_i, x_{i+1}]$ pour $i = 1, \dots, n-1$,
        $S(x)$ est un polynôme cubique (qui peut avoir des coefficients
        différents dans chaque sous-intervalle),
  \item si $g(x)$ est une autre fonction qui satisfait les conditions (a), (b)
        et (c), alors on doit avoir
        \[
          \int_a^b \bigl(S''(x)\bigr)^2\,\dd x \;\le\; \int_a^b \bigl(g''(x)\bigr)^2\,\dd x .
        \]
\end{enumerate}
Dans la terminologie de la mécanique, ces propriétés peuvent s'énoncer de la
manière suivante :
\begin{enumerate}[label=(\alph*)]
  \item la courbe Spline doit passer par tous les n\oe uds ;
  \item la courbe Spline n'est pas brisée ou ne se courbe pas suivant un angle
        pointu ;
  \item entre deux n\oe uds, la courbe Spline est un polynôme de degré 3 ;
  \item la courbe Spline minimise l'énergie potentielle du système. Cette
        propriété implique que
        \[
          S''(a) = S''(b) = 0 .
        \]
\end{enumerate}

\subsubsection{Forme de l'interpolation Spline}

Puisque $S$ est un polynôme cubique pour $x_i \le x \le x_{i+1}$, $S''(x)$ est
alors linéaire et la formule d'interpolation de $S''(x)$ est :
\[
  S''(x) = S''(x_i) + \frac{x - x_i}{x_{i+1} - x_i}\bigl[S''(x_{i+1}) - S''(x_i)\bigr].
\]
Après intégration, on obtient
\[
  S'(x) = S'(x_i) + S''(x_i)(x - x_i)
        + \frac{S''(x_{i+1}) - S''(x_i)}{2(x_{i+1} - x_i)}(x - x_i)^2 .
\]
Une deuxième intégration donne :
\begin{equation}
  S(x) = f_i + S'(x_i)(x - x_i) + \frac{S''(x_i)}{2}(x - x_i)^2
       + \frac{S''(x_{i+1}) - S''(x_i)}{6(x_{i+1} - x_i)}(x - x_i)^3 .
  \tag{a}
\end{equation}
À partir de ces définitions, on peut établir les relations suivantes :
\begin{equation}
  S''_{i-1}h_{i-1} + 2(h_i + h_{i-1})S''_i + S''_{i+1}h_i
  = 6\left[\frac{f_{i+1} - f_i}{h_i} - \frac{f_i - f_{i-1}}{h_{i-1}}\right]
  \qquad (i = 2, \dots, n-1)
  \tag{b}
\end{equation}
où $h_i = x_{i+1} - x_i$, $S''_i = S''(x_i)$.

On a aussi pour $i = 2, \dots, n-1$
\[
  S'_i = \frac{f_{i+1} - f_i}{h_i} - S''_{i+1}\frac{h_i}{6} - S''_i\frac{h_i}{3} .
\]
Les valeurs $S''_1$ et $S''_n$ sont déterminées par la quatrième propriété de
l'interpolation Spline à savoir
\[
  S''(a) = S''(b) = 0 \quad\Longrightarrow\quad S''_1 = S''_n = 0 .
\]
L'équation (b) donne un système algébrique linéaire dont la résolution permet
de trouver tous les $S''_i$. On les substitue dans (a) pour avoir les
fonctions $S(x)$ dans les intervalles $(x_i\,;\,x_{i+1})$.

\subsection{Autres types d'interpolation polynomiale}

Dans les applications, d'autres renseignements peuvent être connus sur la
fonction $f$. Par exemple, si la parité est connue, alors le polynôme est écrit
sous la forme d'une somme de puissances paires de $x$ si la fonction est paire
ou de puissances impaires si la fonction est impaire. Les coefficients du
polynôme peuvent être déterminés comme dans le cas général.

De plus il peut aussi arriver qu'en plus des valeurs $y_k = f(x_k)$, nous
disposons des valeurs des dérivées $y'_k = f'(x_k)$. Dans ces cas on utilise
la formule d'interpolation de \textbf{Hermite}. Il existe également d'autres
types d'interpolation polynomiale telle que celle de \textbf{Stirling} définie
à partir des différences centrées sur base entière, et celle de
\textbf{Bessel} définie à partir des différences centrées sur base demi
entière.

La formule de Stirling s'utilise surtout lorsque $x_p$ est proche d'un n\oe ud
$x_k$ tandis que celle de Bessel est adaptée au cas où $x_p$ est proche de
\[
  x_{k+1/2} = \frac{x_k + x_{k+1}}{2} .
\]

\paragraph{Remarque} L'opérateur de différence centrée sur base entière est
défini par :
\[
  \delta f(x) = f\Bigl(x + \frac{h}{2}\Bigr) - f\Bigl(x - \frac{h}{2}\Bigr)
  \;\longrightarrow\;
  \delta f_k = f_{k+1/2} - f_{k-1/2} .
\]

\section{Interpolation trigonométrique, double et inverse}

\subsection{Interpolation trigonométrique}

Quand la fonction recherchée est périodique, on utilise l'interpolation
trigonométrique. Ce problème a été d'abord analysé par Gauss et ensuite par
Hermite.

La formule d'interpolation de Hermite est la suivante :
\[
  f_k(x) = \prod_{\substack{i=1\\ i\neq k}}^{n+1}\sin(x - x_i)
           \Bigg/ \prod_{\substack{i=1\\ i\neq k}}^{n+1}\sin(x_k - x_i)
  \qquad\text{et}\qquad
  f(x) = \sum_{k=1}^{n+1} f_k(x)\,y_k .
\]
La différence entre la formule de Gauss et celle de Hermite réside dans le
fait que chez Gauss les arguments de la fonction sinus sont divisés par deux.

\subsection{Interpolation double}

\subsubsection{Problème}

Au cours de certaines études, l'on peut être confronté à un problème
d'interpolation à deux arguments, notamment dans le cas des problèmes plans où
une grandeur physique $z$ dépend des coordonnées $x$ et $y$. Ce type de
problème peut se résoudre de deux manières :
\begin{enumerate}[label=(\roman*)]
  \item on peut d'abord interpoler par rapport à une variable, ensuite par
        rapport aux autres ;
  \item on peut aussi trouver une formule d'interpolation qui utilise les
        différences doubles.
\end{enumerate}

\subsubsection{Différences doubles ou à deux voies}

Soit la fonction $Z = z(x,y)$ dont la valeur en un point $(x_m, y_e)$ est
$z_{me}$. On note
\[
  \begin{aligned}
    \Delta^{10} z_{me} &= \Delta_x z_{me} = z_{m+1,e} - z_{m,e}\\
    \Delta^{01} z_{me} &= \Delta_y z_{me} = z_{m,e+1} - z_{me}
  \end{aligned}
\]
Par exemple
\[
  \begin{aligned}
    \Delta^{10} z_{11} &= z_{21} - z_{11}\\
    \Delta^{11} z_{11} &= \Delta^2_{xy} z_{11} = \Delta^{10}(z_{12} - z_{11})
                        = \Delta^{01} z_{21} - \Delta^{01} z_{11} .
  \end{aligned}
\]
De façon générale
\[
  \Delta^{\alpha\beta} z_{11} = \Delta^{\alpha 0} z_{1,\beta+1}
    - \beta\,\Delta^{\alpha 0} z_{1,\beta}
    + \frac{\beta(\beta-1)}{2}\,\Delta^{\alpha 0} z_{1,\beta-1}
    - \dots + (-1)^{\beta}\,\Delta^{\alpha 0} z_{11} .
\]

\subsubsection{Formule d'interpolation double}

Elle est donnée par
\[
  \begin{aligned}
    &z(x,y) = f(x,y) = z_{11}
      + \frac{x - x_1}{h}\Delta^{10} z_{11} + \frac{y - y_1}{l}\Delta^{01} z_{11}\\
    &\quad + \frac{1}{2!}\left[\frac{(x - x_1)(x - x_2)}{h^2}\Delta^{20} z_{11}
        + \frac{2(x - x_1)(y - y_1)}{hl}\Delta^{11} z_{11}
        + \frac{(y - y_1)(y - y_2)}{l^2}\Delta^{02} z_{11}\right]\\
    &\quad + \dots\\
    &\quad + \frac{1}{m!}\Biggl[\frac{(x - x_1)(x - x_2)\cdots(x - x_m)}{h^m}\Delta^{m0} z_{11}
        + \frac{m(x - x_1)\cdots(x - x_{m-1})(y - y_1)}{h^{m-1}l}\Delta^{(m-1)1} z_{11}\\
    &\qquad\quad + \frac{m(m-1)}{2}\,\frac{(x - x_1)\cdots(x - x_{m-2})(y - y_1)(y - y_2)}{h^{m-2}l^2}
                \Delta^{(m-2)2} z_{11} + \dots\\
    &\qquad\quad + \frac{(y - y_1)(y - y_2)\cdots(y - y_m)}{l^m}\Delta^{0m} z_{11}\Biggr]
  \end{aligned}
\]
où $h$ et $l$ sont les pas suivant $x$ et suivant $y$.

\subsection{Interpolation inverse}

Elle consiste à déterminer $x_p$ connaissant $y_p$. On peut dans ce cas
utiliser les formules d'interpolation décrites précédemment. Il suffira de
remplacer $y$ par $x$.

\section{Approximation par la méthode des moindres carrés}

\subsection{Position du problème}

La méthode d'approximation par les moindres carrés est utilisée dans de
nombreuses applications sous des noms différents : optimisation linéaire,
méthode de régression, lissage des courbes.

Les expériences en physique peuvent contenir des milliers de mesures.
L'utilisation des méthodes d'interpolation va alors conduire à des grandes
erreurs d'arrondi et des calculs assez longs. La méthode des moindres carrés
nous permet d'utiliser plusieurs fonctions pour obtenir une approximation
simple.

Par exemple, soit $f_1, f_2, \dots, f_m$ les valeurs correspondantes aux
points $x_1, x_2, \dots, x_m$ et soit
\[
  f(x) = a_0 + a_1x + a_2x^2 + \dots + a_nx^n
\]
un polynôme d'approximation de degré $n$ avec $n \ll m$.

Il est clair que puisque $n \ll m$, il n'est pas possible d'avoir $f$ telle
que $\forall k$, $f(x_k) = f_k$. Mais on peut choisir $f$ telle que
$f(x_k) \approx f_k$. La fonction $f$ sera bonne si la quantité (reste)
\[
  R = \bigl[f_1 - f(x_1)\bigr]^2 + \bigl[f_2 - f(x_2)\bigr]^2 + \dots
      + \bigl[f_m - f(x_m)\bigr]^2
\]
est la plus petite possible. Il s'agit donc de trouver la fonction $f$ qui
minimise le reste $R$.

\subsection{Résolution du problème de l'approximation par les moindres
carrés}

Pour résoudre ce problème, on doit établir un système d'équations linéaires de
la forme :
\[
  \frac{\partial R}{\partial a_k} = 0 \qquad\text{avec } k = 0, \dots, n .
\]
La solution de ce système permet alors de déterminer les coefficients $a_i$
de $f$.

Par exemple, considérons le cas polynomial de premier degré $(n = 1)$ où
\[
  f(x) = a_0 + a_1x .
\]
Le reste est
\[
  R = \sum_{k=1}^{m}\bigl[f_k - (a_0 + a_1x_k)\bigr]^2 .
\]
Pour que $R$ soit minimal, il faut que
$\dfrac{\partial R}{\partial a_0} = 0$ et
$\dfrac{\partial R}{\partial a_1} = 0$.

On trouve alors
\[
  a_0 = \frac{\sum x_k^2\sum f_k - \sum x_k\sum x_kf_k}{m\sum x_k^2 - \left(\sum x_k\right)^2}
  \qquad\text{et}\qquad
  a_1 = \frac{m\sum x_kf_k - \sum x_k\sum f_k}{m\sum x_k^2 - \left(\sum x_k\right)^2}
\]
(établir en exercice).

En général, on choisit $f(x)$ sous la forme
\[
  f(x) = a_0g_0(x) + a_1g_1(x) + \dots + a_ng_n(x)
\]
où les fonctions $g_k(x)$ peuvent être de natures différentes (polynôme,
fonction rationnelle, logarithme, exponentielle, etc.).

Le système d'équations à résoudre est alors défini de la manière suivante :
\[
  \left\{
  \begin{aligned}
    A_{00}a_0 + A_{01}a_1 + \dots + A_{0n}a_n &= b_0\\
    A_{10}a_0 + A_{11}a_1 + \dots + A_{1n}a_n &= b_1\\
    &\;\;\vdots\\
    A_{n0}a_0 + A_{n1}a_1 + \dots + A_{nn}a_n &= b_n
  \end{aligned}
  \right.
\]
où
\[
  b_k = \sum_{i=1}^{m} g_k(x_i)\,f_i, \qquad
  A_{kj} = \sum_{i=1}^{m} g_k(x_i)\,g_j(x_i)
  \qquad\text{avec}\quad
  \left\{
  \begin{aligned}
    j &= 0, 1, \dots, n\\
    k &= 0, 1, \dots, n
  \end{aligned}
  \right.
\]

\begin{exercice}[Exercice 1]
Écrire les algorithmes de calculs des coefficients $b_k$ et $A_{kj}$.
\end{exercice}

\begin{exercice}[Exercice 2]
Les résultats d'une expérience sont donnés dans le tableau suivant :
\begin{center}
  \begin{tabular}{|c|c|c|c|c|c|}
    \hline
    $x$ & 1 & 2 & 3 & 4 & 5\\ \hline
    $y$ & 1,6 & 2 & 2,3 & 2,4 & 2,5\\ \hline
  \end{tabular}
\end{center}
On se propose de déterminer l'expression $y = f(x)$. Pour cela, on choisit
\[
  \text{a)}\quad y = a_0 + \frac{a_1}{x}
  \qquad\text{et}\qquad
  \text{b)}\quad y = a_0 + \frac{a_1}{x} + a_2e^{x} .
\]
En utilisant la méthode des moindres carrés, calculer les coefficients $a_0$,
$a_1$ et $a_2$.
\end{exercice}

\begin{exercice}[Exercice 3]
Du tableau suivant utiliser la formule d'interpolation de Lagrange pour
déterminer $y_p$ pour $x_p = 102$ et $x_p$ pour $y_p = 13{,}5$.
\begin{center}
  \begin{tabular}{|c|c|c|c|c|c|}
    \hline
    $x$ & 93,0 & 96,2 & 100,0 & 104,2 & 108,7\\ \hline
    $y$ & 11,38 & 12,80 & 14,70 & 17,07 & 19,91\\ \hline
  \end{tabular}
\end{center}
\textit{Réponse :} $(x_p = 102$ et $y_p = 15{,}79)$, $(x_p = 97{,}65$ et
$y_p = 13{,}5)$.
\end{exercice}
